\paragraph{哲学。}\label{sub:philosophy}
在任何一类\trCategory{}中构造\trObject{}时，一个广为采用的策略是：先用假想的%
表示该\trObject{}的\trFunctor{}来刻画它，然后再去处理存在性（可表示性）问题。
在我们这里，合适的形式体系是 \trAinfty-模与\trYonedaEmbedding。固定一个%
\trACategory{} $\scr{A}$，并记 $\scr{Q}=\mathit{mod}(\scr{A})$。设 $\scr{M}$ 是一个%
\trAinfty-模。假设存在 $Y\in\Ob\scr{A}$，其在\trYonedaEmbedding{}下的像%
$\scr{Y}$ 在 $H^0(\scr{Q})$ 中与 $\scr{M}$ \trIsomorphic，并在该\trCategory{}中%
取定一个\trIsomorphism{} $[t]\colon\scr{Y}\to\scr{M}$。我们称这一对数据%
$(Y,[t])$ \textbf{\trQuasiRepresent}（quasi-represents）$\scr{M}$。\cref{thm:yoneda-cohomologically-full-faithful}蕴含，
若 $(\tilde{Y},[\tilde{t}])$ 是另一对这样的数据，则在 $H^0(\scr{A})$ 中存在一个%
择定的\trIsomorphism{} $[a]\colon Y\to\tilde{Y}$，满足
$[\tilde{t}]\cdot[\scr{J}^1(a)]=[t]$。滥套术语地说，我们有时会径直把 $Y$ 称作%
\trQuasiRepresent{} $\scr{M}$ 的\trObject，但 $[t]$ 的选取按理仍是数据的一部分。借助%
\trChainMap{} $\lambda=\lambda_{\scr{M}}\colon\scr{M}(Y)
\to\hom_{\scr{Q}}(\scr{Y},\scr{M})$——即\eqref{eq:yoneda-evaluation-chain-map}
——可以给出一个更具体的刻画：

\begin{lemma}\label{thm:quasi-represents-characterization}
$(Y,[t])$ \trQuasiRepresent{} $\scr{M}$，当且仅当下述条件成立：存在 $c\in\scr{M}^0(Y)$，%
满足 $\mu^1_{\scr{M}}(c)=0$ 且 $[t]=[\lambda(c)]$，使得对一切 $X\in\Ob\scr{A}$，映射
\begin{equation}\label{eq:quasi-representing-map}
\hom_{\scr{A}}(X,Y)\longrightarrow\scr{M}(X),\quad
b\longmapsto(-1)^{|b|}\mu^2_{\scr{M}}(c,b)
\end{equation}
是\trQuasiisomorphism。
\end{lemma}

\begin{proof}
由于 $\mu^1_{\scr{M}}(c)=0$，$\lambda(c)$ 是一个从 $\scr{Y}$ 到 $\scr{M}$ 的%
\trAinfty-模\trHomomorphism，其一阶项除正负号外正是
\eqref{eq:quasi-representing-map}。\cref{thm:module-composition-quasiisomorphism}
蕴含 $[t]=[\lambda(c)]$ 是 $H^0(\scr{Q})$ 中的\trIsomorphism，故 $(Y,[t])$ \trQuasiRepresent%
$\scr{M}$。反之，若有一对\trQuasiRepresent{} $\scr{M}$ 的数据 $(Y,[t])$，则由\cref{thm:yoneda-evaluation-quasiisomorphism}，%
存在 $c$ 使得 $\mu^1_{\scr{M}}(c)=0$ 且 $[t]=[\lambda(c)]$。由于 $[\lambda(c)]$
是 $H^0(\scr{Q})$ 中的\trIsomorphism，它特别地在底层\trCohomology{}模之间也%
诱导\trIsomorphism，而这蕴含\eqref{eq:quasi-representing-map} 对一切 $X$ 都是%
\trQuasiisomorphism。
\end{proof}
